% Original English questions and concise English solutions. Compile using XeLaTeX.
\documentclass[UTF8,10pt,letterpaper,fontset=fandol]{ctexart}
\usepackage[margin=1in]{geometry}
\usepackage{amsmath,amssymb}
\usepackage{enumitem}
\usepackage{needspace}
\usepackage[unicode,hidelinks]{hyperref}
\pagestyle{plain}
\raggedbottom
\setlength{\parindent}{0pt}
\setlength{\parskip}{1pt}
\setlength{\emergencystretch}{2em}
\setlist[enumerate]{label=(\arabic*),leftmargin=2em,itemsep=0.3em,topsep=0.2em}
\setcounter{secnumdepth}{0}
\ctexset{section={format=\Large\bfseries,beforeskip=1.2ex,afterskip=0.8ex}}
\newcommand{\F}{\mathbb F}
\newcommand{\R}{\mathbb R}
\newcommand{\C}{\mathbb C}
\newcommand{\Ideal}[1]{\langle #1\rangle}
\newcommand{\Solution}{\par\Needspace{4\baselineskip}\textbf{Solution}\par}
\begin{document}
\begin{center}
{\LARGE Homework 3}\par
\vspace{0.6em}
October 15, 2026\par
Name: YuanHao Tang\quad Student ID: 3120265621
\end{center}

\section{Problem 1}
Let \(C_{n-1}(t)\) be the space of continuous piecewise linear functions on \([0,n-1]\)
whose slope is constant on each interval \([i-1,i]\), \(i=1,\ldots,n-1\).
Define \(\mathcal A:C_{n-1}(t)\to C_{n-1}(t)\) by
\[
\mathcal A[f(t)]=f(n-1-t).
\]
\begin{enumerate}
\item Prove that \(\mathcal A\) is a linear transformation.
\item Let \(\epsilon_j(t)\in C_{n-1}(t)\), \(j=1,\ldots,n\), be the nodal basis satisfying
\(\epsilon_j(i-1)=\delta_{ij}\) for \(i,j=1,\ldots,n\).
Determine the matrix representation of \(\mathcal A\) with respect to this basis.
\end{enumerate}

\Solution
\begin{enumerate}
\item Reflection \(t\mapsto n-1-t\) maps grid intervals to grid intervals, preserving continuity
and piecewise linearity. Thus the map is well-defined. For functions \(f,g\) and scalars \(a,b\),
\[
\begin{aligned}
\mathcal A(af+bg)(t)&=a f(n-1-t)+b g(n-1-t)\\
&=a(\mathcal Af)(t)+b(\mathcal Ag)(t),
\end{aligned}
\]
so \(\mathcal A\) is linear.
\item At every node,
\[
(\mathcal A\epsilon_j)(i-1)=\epsilon_j(n-i)=\delta_{n+1-i,j}=\delta_{i,n+1-j}.
\]
Since nodal values uniquely determine a function, \(\mathcal A\epsilon_j=\epsilon_{n+1-j}\). Hence
\[
\boxed{[\mathcal A]_{\epsilon}=J_n=(\delta_{i,n+1-j})_{i,j=1}^n},\qquad
(J_n)_{ij}=\begin{cases}1,&i+j=n+1,\\0,&i+j\ne n+1.\end{cases}
\]
This is the reversal matrix, with ones on the antidiagonal and zeros elsewhere.
\end{enumerate}

\section{Problem 2}
Let \(I_1\) and \(I_2\) be the polynomial ideals generated by
\[
a_1(\lambda)=(\lambda+1)(\lambda+2)^2,\qquad
a_2(\lambda)=(\lambda+2)(\lambda+3)(\lambda+5),
\]
respectively. Compute \(I_1+I_2\) and \(I_1\cap I_2\).

\Solution
Work in \(\F[\lambda]\), where \(\F=\R\) or \(\C\). Using
\(I_1+I_2=\Ideal{\gcd(a_1,a_2)}\) and
\(I_1\cap I_2=\Ideal{\operatorname{lcm}(a_1,a_2)}\), take the smaller and larger
multiplicities of each factor, respectively. Thus
\[
\boxed{I_1+I_2=\Ideal{\lambda+2}},
\]
\[
\boxed{I_1\cap I_2=\Ideal{(\lambda+1)(\lambda+2)^2(\lambda+3)(\lambda+5)}}.
\]

\section{Problem 3}
Let \(I_1\) and \(I_2\) be the polynomial ideals generated by
\[
\begin{aligned}
a_1(\lambda)&=\lambda^4+\lambda^3-\lambda-1,\\
a_2(\lambda)&=\lambda^4+2\lambda^3+3\lambda^2+2\lambda+1,
\end{aligned}
\]
respectively. Compute \(I_1+I_2\) and \(I_1\cap I_2\).

\Solution
In the same coefficient field as Problem 2, set \(q(\lambda)=\lambda^2+\lambda+1\). Factoring gives
\[
a_1=(\lambda+1)(\lambda-1)q=(\lambda^2-1)q,\qquad a_2=q^2.
\]
Since \(q(1)=3\) and \(q(-1)=1\), we have \(\gcd(\lambda^2-1,q)=1\). Therefore
\[
\boxed{I_1+I_2=\Ideal{\lambda^2+\lambda+1}},
\]
\[
\boxed{I_1\cap I_2=\Ideal{(\lambda^2-1)(\lambda^2+\lambda+1)^2}}.
\]
\end{document}
